\documentclass[11pt,a4paper]{article}
\usepackage{amsmath,amssymb,amsthm}
\usepackage[margin=26mm]{geometry}
\usepackage[hidelinks]{hyperref}
\newtheorem{theorem}{Theorem}
\newtheorem{lemma}[theorem]{Lemma}
\newtheorem{proposition}[theorem]{Proposition}
\newtheorem{corollary}[theorem]{Corollary}
\theoremstyle{definition}
\newtheorem*{conjecture}{Conjecture}

% Frame statements while keeping proofs in the ordinary text flow.
\newsavebox{\statementbox}
\newenvironment{statementframe}
  {\par\addvspace{4pt}\noindent
   \setlength{\fboxsep}{6pt}\setlength{\fboxrule}{0.45pt}%
   \begin{lrbox}{\statementbox}%
   \begin{minipage}{\dimexpr\linewidth-2\fboxsep-2\fboxrule\relax}}
  {\end{minipage}\end{lrbox}%
   \fbox{\usebox{\statementbox}}\par\addvspace{4pt}}
\newcommand{\R}{\mathbb R}
\newcommand{\C}{\mathbb C}
\newcommand{\1}{\mathbf1}
\newcommand{\A}{\mathcal A}
\newcommand{\F}{\mathcal F}
\newcommand{\E}{\mathbb E}
\newcommand{\norm}[1]{\left\lVert #1\right\rVert}
\setlength{\parindent}{1.2em}
\setlength{\parskip}{3pt}
\setlength{\emergencystretch}{2em}
\title{Pointwise Convergence for Two\\ Commuting Transformations\\[5pt]
\large Revised proof manuscript with traceable inputs}
\author{}
\date{October 8, 2026}
\begin{document}
\maketitle
\vspace{-2.4em}
\begin{abstract}
We present a proof argument for pointwise convergence of bilinear
averages associated with two commuting invertible measure-preserving
transformations. Finite-matrix and heat-flow lemmas yield a smooth
kernel estimate through higher Gaussian derivatives and an entrywise
switching argument. A planar maximal estimate, spatial transference,
closure in the kernel's $L^1$ norm, and exact phase reconstruction then
give convergence for bounded functions. Truncation extends the
conclusion to $L^2\times L^2$. This is a proof manuscript for review;
the qualitative main theorem now has a compiled Lean formalization.
The formal statements, their standard axiom audit, and the scope of the
formalization are recorded in Appendix~\ref{app:formal}.
Independent mathematical review and the separate full dependency-source
rebuild remain distinct validation steps.
\end{abstract}

\section{Main statement and planar estimate}

Let $(X,\mathcal B,\mu)$ be a probability space. Let $T,S$ be invertible,
measurable, measure-preserving transformations with measurable inverses
and $TS=ST$. For measurable complex-valued functions $f,g$ and $N\ge1$, put
\[
 A_N(f,g)(x)=\frac1N\sum_{n=0}^{N-1} f(T^nx)g(S^nx).
\]
\begin{statementframe}
\begin{conjecture}
For every such system and every $f,g\in L^\infty(\mu)$, the sequence
$A_N(f,g)(x)$ converges for $\mu$-almost every $x$.
\end{conjecture}
\end{statementframe}
The framed conjecture records the question being addressed.
Theorem~\ref{thm:main} is the conclusion proved below.
All convergence assertions are almost everywhere, for each fixed pair of
inputs; they do not assert convergence at every point or simultaneously
for every equivalence class.

Set
\[
 \varphi(t)=(2\pi)^{-1/2}e^{-t^2/2},\qquad
 E(t)=(3-t^2)\varphi(t),\qquad O(t)=(3t-t^3)\varphi(t),
\]
and write $D_rK(t)=r^{-1}K(t/r)$ for $r>0$. For $K\in L^1(\R)$ and
bounded measurable $F,G:\R^2\to\C$, define
\[
 \A_N^K(F,G)(u,v)=\int_\R K(t)F(u+Nt,v)G(u,v+Nt)\,dt.
\]
An $m$-chain is a tuple $\boldsymbol n=(n_0,\ldots,n_m)$ of positive
integers with $n_0<\cdots<n_m$. For a sequence $a=(a_N)$ set
\[
 V_{\boldsymbol n}(a)=\sum_{j=1}^m |a_{n_j}-a_{n_{j-1}}|,
 \qquad M_{\F}(a)=\max_{\boldsymbol n\in\F}V_{\boldsymbol n}(a),
\]
where $\F$ is a finite nonempty family of $m$-chains. Let
$Q_L=[-L,L]^2$, with $|Q_L|=4L^2$.

\begin{statementframe}
\begin{proposition}[Planar finite maximal estimate]\label{prop:planar}
For each $K\in\{D_rE,D_rO:r>0\}$ and $C_f,C_g\ge0$, there are finite
numbers $C_m\ge0$, $m\ge1$, such that $C_m/m\to0$ and
\begin{equation}\label{eq:H}
 \int_{\R^2} M_{\F}\bigl((\A_N^K(F,G)(p))_{N\ge1}\bigr)\,dp
 \le |Q_L|C_m
\end{equation}
for every $L>0$, every bounded measurable $F,G$ vanishing outside $Q_L$
with $|F|\le C_f$, $|G|\le C_g$, and every finite nonempty family $\F$
of $m$-chains. The sequence $(C_m)$ is independent of $L,F,G,\F$.
\end{proposition}
\end{statementframe}
\begin{proof}
Proposition~\ref{prop:smooth} establishes the smooth kernel estimate
from the finite-matrix and heat-flow lemmas stated in
Section~\ref{sec:analytic}. Lemma~\ref{lem:smooth-to-H} converts that
estimate into \eqref{eq:H}, with $C_m=C_\kappa C_fC_g\sqrt m$ for
$K=D_r\kappa$, $\kappa=E,O$. In particular $C_m/m\to0$.
\end{proof}

The maximum in \eqref{eq:H} is inside the integral. The proof therefore
requires an estimate uniform over output-dependent choices of chains,
as supplied by the switching argument in Section~\ref{sec:analytic}.

\begin{statementframe}
\begin{theorem}\label{thm:main}
For every probability space and pair $T,S$ as above,
and every $f,g\in L^2(\mu)$, the sequence $A_N(f,g)$ converges almost
everywhere. In particular, the conclusion holds for
$f,g\in L^\infty(\mu)$ and for $f,g\in L^3(\mu)$.
\end{theorem}
\end{statementframe}

We prove this theorem below. Throughout, representatives may be chosen
measurable; a bounded equivalence class may be represented by a function
bounded everywhere. Changing representatives does not affect the discrete
averages outside the countable union of the inverse images of their
exceptional sets under $T^n$ and $S^n$.

\section{The analytic estimate}\label{sec:analytic}

We use the finite-matrix and heat-flow lemmas stated below, with
their proofs in the frozen source of \cite{FrozenMatrix,FrozenHeat}.
The precise source labels are listed in Appendix~\ref{app:inputs}.
These are proof manuscripts, and their citation does not denote external
peer review. Their fixed commit and locally recorded source hashes make
the inputs traceable independently of the current state of the repository.
We give the higher-derivative and switching argument that derives the
planar estimate from these lemmas. The status of the mathematical
review and the status of Lean formalization are separate from the
logical dependencies of this written argument.

Write $g_\sigma(t)=(2\pi\sigma)^{-1/2}e^{-t^2/(2\sigma)}$.
For finite arrays $A_v$, with indices $v$ taken modulo $3$, use a mesh
$h>0$, nodes $u_i=ih$, and weights
\[
 w_v(i)=h g_{\sigma_v}(u_i-p_v),\quad
 W_v(i,j)=A_v(i,j)\sqrt{w_v(i)w_{v+1}(j)},\quad
 n_v=\left(\sum_{i,j}h^2|A_v(i,j)|^3\right)^{1/3}.
\]
Put $R_v=[W_v\ W_{v-1}^*]$, $T_v=R_v^*R_v$, $S_v=R_vR_v^*$, and
\[
 \Phi(P)=\operatorname{tr}(P^{3/2}),\qquad
 b_P(K)=\frac32\sum_{i,j}\frac{|K_{ij}|^2}{\sqrt{p_i}+\sqrt{p_j}}.
\]
Here the sum is taken in an eigenbasis of the positive support of $P$,
and $K$ is Hermitian and supported there. We retain the following
matrix inputs: joint convexity of $(P,K)\mapsto b_P(K)$ under probability
averages with a fixed support and integrable operator norms; and
\begin{equation}\label{eq:matrix-input}
 |\operatorname{tr}(D_0W_0D_1W_1D_2W_2)|
 \le\frac53\max_v\|D_v\|_{\mathrm{op}}
       \sum_v b_{T_v}(R_v^*D_vR_v)
\end{equation}
for Hermitian $D_v$. We also retain the rank-change differentiation
formula
\[
 \Phi(R_+^*R_+)-\Phi(R_-^*R_-)
 =3\int_0^1\operatorname{Re}\langle R_z\sqrt{R_z^*R_z},R_+-R_-\rangle_{\mathrm{HS}}\,dz
\]
with $R_z=R_-+z(R_+-R_-)$, and the entry bound
\[
 |(R\sqrt{R^*R})_{ij}|
 \le\|\operatorname{row}_iR\|_2\|\operatorname{col}_jR\|_2.
\]

Let $C_{v,w}$ denote the block of $R_v$ belonging to its neighbor $w$.
Set
\[
 N_v=\operatorname{diag}_i\frac{u_i-p_v}{\sigma_v},\quad
 U_v=R_v^*N_vR_v,\quad V_v=b_{T_v}(U_v),
\]
\[
 P_{v,w}=C_{v,w}N_wC_{v,w}^*,\qquad
 Y_{v,w}=b_{S_v}(P_{v,w}).
\]
The polarization of $b_{S_v}$ applied to the two column insertions is
 denoted $Z_v$; in particular $2Z_v\le\sum_{w\ne v}Y_{v,w}$.
Let $\Pi_d=\{p_0+p_1+p_2=d\}$, with any two centers as free Lebesgue
coordinates, and $J_v(d,s)=\int_{\Pi_d}\Phi(T_v)\,d\pi_d$, where
$\sigma_i=\alpha_i s$. The heat-flow inputs are
\begin{align}\label{eq:heat-input}
 2\partial_s J_v
 &=\Bigl(\sum_i\alpha_i\Bigr)\partial_d^2J_v
   -\int_{\Pi_d}\Bigl(\alpha_vV_v+\sum_{w\ne v}\alpha_wY_{v,w}\Bigr),\notag\\
 \partial_d^2J_v&=\int_{\Pi_d}Z_v.
\end{align}
We retain the initial bound $J_v\le C\sum_i n_i^3$ whenever all
$\sigma_i\ge h^2$. For a single retained edge, its energy $J^e_{v,w}$
is independent of $d$ and obeys
\[
 -2\partial_sJ^e_{v,w}
 =\iint(\alpha_vV^e_{v,w}+\alpha_wY^e_{v,w})\,dp_v\,dp_w,
 \qquad Y_{v,w}\le Y^e_{v,w}.
\]
For an edge with discrete $L^3$ norm $n_e$, the initial bound is
\[
 J^e_{v,w}(s_0)\le M n_e^3,\qquad M=1+\sqrt{2/\pi},
\]
provided both endpoint variances are at least $h^2$; see
\cite[Lemmas \texttt{lem:initial-energy} and
\texttt{lem:single-edge-energy}]{FrozenHeat}.
The absolute factor $M$ is absorbed in $C$ in the subsequent budget.
Gaussian domination of the needed center derivatives on compact positive
scale intervals, and the determinant-one changes between plane charts,
are justified by the proof of
\cite[Lemma \texttt{lem:heat-identities}]{FrozenHeat}.

\begin{statementframe}
\begin{proposition}[Smooth kernel estimate]\label{prop:smooth}
For bounded measurable compactly supported $F,G$ and $\kappa=E,O$,
\begin{equation}\label{eq:smooth}
 \|\mathcal C_m(B^\kappa(F,G))\|_{3/2}
 \le C_\kappa\sqrt m\,\|F\|_3\|G\|_3,
\end{equation}
where
\[
 B_L^\kappa(F,G)(x,y)=\int_\R F(x+t,y)G(x,y+t)D_L\kappa(t)\,dt,
\]
and $\mathcal C_m$ is the pointwise supremum of sums of increments over
at most $m$ ordered, nonoverlapping positive scale intervals. The bound
is uniform over finite menus and output-dependent choices.
\end{proposition}
\end{statementframe}
\begin{proof}
Put $s=L^2$ and $k_{r,L}=L^r g_{s/3}^{(r)}$. The ratio
$k_{r,L}/g_s$ is bounded uniformly in $L$, for each fixed $r$.
Narrow the row variance at vertex $v$ to $s/6$, move its center by $d$,
and keep its column weights fixed. Explicitly, if the base centers
$q$ lie in $\Pi_0$, put $p_v=q_v+d$ and $p_w=q_w$ for $w\ne v$.
Then $p\in\Pi_d$, and the map between the two planes has absolute
Jacobian one in the two free-center coordinates. The variance vector
for this fixed vertex is
$\alpha^{(v)}_v=1/6$ and $\alpha^{(v)}_w=1$ for $w\ne v$.
Denote the new Gram and score insertion by $T_v^d,U_v^d$. With $d\nu_s(d)=g_{5s/6}(d)\,dd$, Gaussian
convolution gives
\[
 T_v^{\mathrm{base}}=\int T_v^d\,d\nu_s,
 \qquad R_v^*\operatorname{diag}(k_{r,L}/g_s)R_v
       =\int c_{r,s}(d)U_v^d\,d\nu_s,
\]
\[
 c_{r,s}(d)=-\frac{L^r g_{s/6}^{(r-1)}(d)}{g_{5s/6}(d)},
 \qquad |c_{r,s}(d)|\le C_r\sqrt s.
\]
All row weights are positive, so the Gram support is independent of
$d$. In finite dimensions the relevant operator norms are integrable
against $\nu_s$. Joint convexity and quadratic homogeneity therefore give
\begin{equation}\label{eq:score}
 b_{T_v^{\mathrm{base}}}\bigl(R_v^*\operatorname{diag}(k_{r,L}/g_s)R_v\bigr)
 \le C_rs\int V_v^d\,d\nu_s.
\end{equation}
The signed coefficient remains in the insertion; it is not used as a
probability density.

For this narrowed vertex write $J_v^{(v)}(d,s)$ for the plane energy
with variance vector $\alpha^{(v)}$ just specified, and define
$\widetilde J_v=\int J_v^{(v)}(d,s)\,d\nu_s(d)$.
Use a tilde for the corresponding averaged costs. Thus the symbol
$J_v^{(v)}$ here does not denote the base equal-variance energy. The variance rates
sum to $13/6$. Differentiating the averaging density adds heat rate
$5/6$. Thus \eqref{eq:heat-input} and integration by parts give
\[
 2\partial_s\widetilde J_v
 =3\widetilde Z_v-\tfrac16\widetilde V_v-\sum_{w\ne v}\widetilde Y_{v,w},
\]
whence
\begin{equation}\label{eq:averaged-energy}
 \tfrac16\widetilde V_v
 \le-2\partial_s\widetilde J_v+\tfrac12\sum_{w\ne v}\widetilde Y_{v,w}.
\end{equation}
The last terms are controlled by single-edge dissipation: their plane
integrals do not depend on $d$, and each column rate is $1$.

Now change only edge $0$, using at most $m$ nonoverlapping scale
intervals and coefficients of modulus at most $1$ at each $(i,j)$:
\[
 B_0^s(i,j)=A_0(i,j)\sum_\ell\beta_\ell(i,j)
                  \1_{\{a_\ell(i,j)^2<s<b_\ell(i,j)^2\}}.
\]
Then $|B_0^s|\le|A_0|$ and
$\sum_\tau|\Delta B_0^\tau(i,j)|\le2m|A_0(i,j)|$.
At a switch, use the original arrays to define
\begin{align*}
 X^2&=\sum_l|A_0(i,l)|^2w_1(l),&Y^2&=\sum_l|A_0(l,j)|^2w_0(l),\\
 H^2&=\sum_l|A_2(l,i)|^2w_2(l),&Q^2&=\sum_l|A_1(j,l)|^2w_2(l).
\end{align*}
The rank-change formula and the row-column bound yield
\[
 |\Delta e_0|+|\Delta e_1|
 \le C\sum_{i,j}w_0(i)w_1(j)|\Delta B_0(i,j)|(2XY+HY+QX),
\]
where $e_v=\Phi(T_v)$. Star $2$ does not change.
For fixed $i,j$, use the probability measure
\[
 h^{-2}w_0(i)w_1(j)\,d\pi_d\,d\nu_s(d).
\]
Gaussian convolution gives
\begin{align*}
 \E w_0(l)&=h g_{2\alpha_0s}(u_l-u_i),&
 \E w_1(l)&=h g_{2\alpha_1s}(u_l-u_j),\\
 \E w_2(l)&=h g_{3s}(u_l+u_i+u_j).
\end{align*}
Gaussian averages are bounded by the discrete one-dimensional maximal
operator, uniformly when the variances exceed the mesh scale; the
Gaussian domination and strong $L^{3/2}$ bound are proved in
\cite[Lemma \texttt{lem:line-maximal}]{FrozenPreliminaries}. Hence
$\E X^2,\E Y^2,\E H^2,\E Q^2$ have scale-independent majorants whose
square roots have discrete $L^3$ norms at most
$Cn_0,Cn_0,Cn_2,Cn_1$, respectively. This follows by the strong
$L^{3/2}$ maximal inequality on each slice; the index $-i-j$ is a
bijection along the relevant slice.

Cauchy--Schwarz for the center probability measure, the total entrywise
switch bound, and discrete H\"older now give
\[
 \sum_\tau\sum_v|\Delta\widetilde J_v|
 \le Cm\{n_0^3+n_0^2(n_1+n_2)\}.
\]
For the varying single edge the same argument costs $Cm n_0^3$.
Fixed edges have no jumps and are not initialized anew at each switch.
Integrating \eqref{eq:averaged-energy} between switches and telescoping
the energies therefore yields
\begin{equation}\label{eq:budget-candidate}
 \int\widetilde V_v(s)\,ds
 \le C\left[\sum_{j=0}^2 n_j^3+
                 m\{n_0^3+n_0^2(n_1+n_2)\}\right].
\end{equation}
All scales lie in a fixed compact interval during this step, and the
mesh is chosen with $h^2\le s_-/6$.

For $r\ge3$, insert the three derivative windows with orders
$(1,1,r-2)$ into \eqref{eq:matrix-input}. Apply \eqref{eq:score} and
integrate over the center plane. Their convolution is
\[
 k_{1,L}*k_{1,L}*k_{r-2,L}=L^r g_s^{(r)}=D_L\varphi^{(r)}.
\]
Integration with $ds/s=2\,dL/L$ and \eqref{eq:budget-candidate}
bounds the trilinear form for output-dependent sums of kernels
\[
 I_{a,b}^{(r)}=\int_a^b D_L\varphi^{(r)}\,\frac{dL}{L}.
\]
Normalize the three input norms to $(m^{-1/2},1,1)$. The bracket in
\eqref{eq:budget-candidate} becomes
$m^{-3/2}+2+m(m^{-3/2}+2m^{-1})\le6$.
Restoring the original norms by trilinearity gives the bound
$C_r\sqrt m\prod_v n_v$.

Fix the actual finite scale menu before any mesh or spatial selector.
Its finitely many positive lower endpoints have a common lower bound
$\ell>0$. The interval-masked scale kernel is identically zero for
$0<L\le\ell$, so removing an auxiliary lower cutoff is eventually an
exact equality. Choose the mesh sufficiently fine relative to this fixed
$\ell$; this threshold is uniform over all selections from the finite menu.
There is no assertion of a common positive lower endpoint over the later
countable collection of all menus.

Fix a rectangular step choice before refining the mesh. Riemann sums then yield the corresponding continuous
trilinear estimate. To pass to measurable choices, keep the valid
interval lists as finitely many labels and approximate their measurable
partition by rectangular partitions in the finite measure having density
the sum of the absolute trilinear integrands. Every label preserves the
switch bound. For each fixed menu all kernels are integrable; H\"older
also gives continuity in the three $L^3$ inputs, so density extends the
estimate to those inputs. Finite menus of rational endpoints and a
countable dense set of phases are then increased monotonically.
More explicitly, for a fixed finite kernel family $(K_q)$ and inputs
$G_0,G_1,G_2\in L^3$, the approximation measure has density
\[
 w(a,b)=\sum_q\int_{\R}
 |G_0(a,b)G_1(b,c)G_2(c,a)K_q(a+b+c)|\,dc.
\]
The substitution $t=a+b+c$ and H\"older on its two free coordinates give
$\int w\le\sum_q\|K_q\|_1\prod_v\|G_v\|_3<\infty$.
The order is therefore: the fixed finite menu and lower-cutoff removal,
the grid limit for rectangle selectors, approximation of the measurable
selector in $w\,da\,db$, $L^3$ input density, and finally directed finite
maxima increasing to the countable supremum inside the integral.
Approximation constants may depend on the fixed kernel family; the final
sharp estimate is independent of it. For the qualitative natural-average
theorem it suffices to use positive natural chains and the four phases
$1,-1,i,-i$, at a fixed harmless factor in the bound. No limit to arbitrary
phases or uncountable endpoints is needed for that formal route.

Finally, direct differentiation gives
\[
 \partial_{\log L}D_LE=-D_L\varphi^{(4)},\qquad
 \partial_{\log L}D_LO=D_L(\varphi^{(5)}+3\varphi^{(3)}).
\]
Use the estimates of orders $3,4,5$, linearize each absolute increment
by a complex phase, and use $L^{3/2}$--$L^3$ duality. The substitution
$x=-a,y=-b,t=a+b+c$ identifies the triangular form with $B_L^\kappa$
and preserves the input norms. This proves \eqref{eq:smooth} from the
stated finite-matrix and heat-flow lemmas.
\end{proof}

\begin{statementframe}
\begin{lemma}\label{lem:smooth-to-H}
The smooth estimate \eqref{eq:smooth} implies the planar estimate
\eqref{eq:H} with $C_m=C_\kappa C_fC_g\sqrt m$ for
$K=D_r\kappa$, $\kappa=E,O$.
\end{lemma}
\end{statementframe}
\begin{proof}
For every $N$, changing variables gives
$\A_N^{D_r\kappa}(F,G)=B_{rN}^\kappa(F,G)$.
Every finite chain maximum is bounded by $\mathcal C_m(B^\kappa)$.
If $F,G$ vanish outside $Q_L$, these averages vanish outside $Q_L$:
$F(u+t,v)$ forces $|v|\le L$, and $G(u,v+t)$ forces $|u|\le L$.
H\"older therefore gives
\begin{align*}
 \int_{\R^2}M_\F(\A^{D_r\kappa})
 &\le |Q_L|^{1/3}\|\mathcal C_m(B^\kappa)\|_{3/2}\\
 &\le |Q_L|^{1/3}C_\kappa\sqrt m\|F\|_3\|G\|_3
 \le |Q_L|C_\kappa C_fC_g\sqrt m.
\end{align*}
The constants do not depend on the finite family or on $L$.
Since $\sqrt m/m\to0$, the constants have the required sublinear growth.
\end{proof}

\section{A convergence criterion and spatial transference}

\begin{statementframe}
\begin{lemma}\label{lem:jump}
Let $(Y,\nu)$ be a probability space and $a_N:Y\to\C$ measurable. Suppose
that $\int_Y M_{\F}(a)\,d\nu\le C_m$ for every finite family of
$m$-chains, where $C_m<\infty$ and $C_m/m\to0$. Then $a_N$ converges
almost everywhere.
\end{lemma}
\end{statementframe}
\begin{proof}
The set of $m$-chains is countable. Enumerating it and applying monotone
convergence to finite maxima gives
\[
 J_m(y):=\sup_{\boldsymbol n}V_{\boldsymbol n}(a(y)),\qquad
 \int_Y J_m\,d\nu\le C_m.
\]
If a sequence is not Cauchy, there is $\eta>0$ such that arbitrarily far
out two terms have distance at least $2\eta$. From any chosen term, at
least one of these two later terms has distance at least $\eta$.
Induction therefore constructs chains of every length with all successive
increments at least $\eta$. Thus $J_m\ge m\eta$ for every $m$.
For each fixed rational $\eta>0$, Markov's inequality gives
\[
 \nu\left(\bigcap_{m\ge1}\{J_m\ge m\eta\}\right)
 \le \inf_{m\ge1}\frac{C_m}{m\eta}=0.
\]
Taking a countable union in $\eta$ proves the Cauchy property almost
everywhere. Completeness of $\C$ gives convergence.
\end{proof}

Let $U_t,V_t$ be commuting, jointly measurable, measure-preserving
$\R$-actions on a probability space $(Y,\nu)$. For bounded measurable
$f,g:Y\to\C$, put
\begin{equation}\label{eq:flowaverage}
 a_N^K(y)=\int_\R K(t)f(U_{Nt}y)g(V_{Nt}y)\,dt.
\end{equation}
The integral is absolutely convergent and defines a measurable function.

\begin{statementframe}
\begin{proposition}\label{prop:transfer}
Fix $K\in L^1(\R)$ and bounds $|f|\le C_f$, $|g|\le C_g$.
If \eqref{eq:H} holds for this kernel, these bounds, and a fixed $m$
with constant $C_m$, then
\[
 \int_Y M_{\F}\bigl((a_N^K(y))_{N\ge1}\bigr)\,d\nu(y)\le C_m
\]
for every finite nonempty family $\F$ of $m$-chains.
\end{proposition}
\end{statementframe}
\begin{proof}
Let $B$ be the largest endpoint in $\F$. For $y\in Y$ define
\[
 F_y(u,v)=f(U_uV_vy),\qquad G_y(u,v)=g(U_uV_vy).
\]
Commutativity gives the exact identity
\[
 \A_N^K(F_y,G_y)(p)=a_N^K(U_{p_1}V_{p_2}y).
\]
Fix $R,L>0$ and replace $F_y,G_y$ by their restrictions to $Q_{L+BR}$,
extended by zero. Denote the resulting functions by $F'_y,G'_y$.
For $p\in Q_L$, $N\le B$, and $|t|\le R$, both shifted points in the
integrand remain in $Q_{L+BR}$. Consequently
\[
 |\A_N^K(F_y,G_y)(p)-\A_N^K(F'_y,G'_y)(p)|
 \le \varepsilon_R,
 \qquad \varepsilon_R=2C_fC_g\int_{|t|>R}|K(t)|\,dt.
\]
Each chain sum changes by at most $2m\varepsilon_R$. The same bound
holds for the finite maximum, independently of the size of $\F$.
Applying \eqref{eq:H} to $F'_y,G'_y$ yields
\[
 \int_{Q_L}M_{\F}\bigl(a^K(U_{p_1}V_{p_2}y)\bigr)\,dp
 \le |Q_{L+BR}|C_m+2m\varepsilon_R|Q_L|.
\]
Integrate in $y$. Tonelli's theorem and measure preservation identify
the left side with $|Q_L|\int_Y M_{\F}(a^K)\,d\nu$. Division by
$|Q_L|>0$ gives
\[
 \int_Y M_{\F}(a^K)\,d\nu
 \le \left(1+\frac{BR}{L}\right)^2C_m+2m\varepsilon_R.
\]
First let $L\to\infty$ with $R$ fixed, then let $R\to\infty$.
Integrability of $K$ proves the claim.
\end{proof}

\begin{statementframe}
\begin{corollary}\label{cor:euler}
The averages \eqref{eq:flowaverage} converge almost everywhere
for each $K=D_rE,D_rO$, $r>0$, and bounded measurable $f,g$.
\end{corollary}
\end{statementframe}
\begin{proof}
Proposition~\ref{prop:planar} supplies the finite maximal estimate.
Apply Proposition~\ref{prop:transfer} for every $m$, followed by
Lemma~\ref{lem:jump}.
\end{proof}

\section{Passage from Gaussian kernels to the interval kernel}

\begin{statementframe}
\begin{lemma}\label{lem:closure}
Let $H_N(y,t)$ be measurable complex-valued functions on $Y\times\R$
with $|H_N(y,t)|\le M$ for every $N,y,t$. The real kernels $K\in L^1(\R)$
for which $\int_\R K(t)H_N(y,t)\,dt$ converges for almost every $y$
form a closed linear subspace $\mathcal C$ of $L^1(\R;\R)$.
\end{lemma}
\end{statementframe}
\begin{proof}
Linearity follows by intersecting finitely many sets of full measure.
Suppose $K_j\in\mathcal C$ and $\norm{K_j-K}_1\to0$. Intersect the
countably many sets of convergence for $K_j$. At any remaining point,
\[
 \sup_N\left|\int_\R(K-K_j)(t)H_N(y,t)\,dt\right|
 \le M\norm{K-K_j}_1.
\]
The Cauchy criterion, first in $N$ and then in $j$, proves convergence
for $K$. This proves closedness.
\end{proof}

\begin{statementframe}
\begin{lemma}\label{lem:gaussian}
If a closed linear subspace $\mathcal C\subset L^1(\R;\R)$ contains
$D_rE$ and $D_rO$ for every $r>0$, it contains $t^de^{-at^2}$ for
every integer $d\ge0$ and every $a>0$.
\end{lemma}
\end{statementframe}
\begin{proof}
For $k=0,1$ put
\[
 E_{k,a}(t)=(3-2at^2)t^ke^{-at^2},\qquad G_{k,a}(t)=t^ke^{-at^2}.
\]
Each $E_{k,a}$ is a constant multiple of a positive dilation of $E$ or
$O$, hence lies in $\mathcal C$. The map $s\mapsto E_{k,a}(s\,\cdot)$
is continuous in $L^1$ on every compact subinterval of $(0,\infty)$:
on $[\delta,1]$ a common integrable bound is a constant times
$(|t|^k+|t|^{k+2})e^{-a\delta^2t^2}$, so dominated convergence applies.
For $0<\delta<1$, integration of the derivative of
$s^3t^ke^{-as^2t^2}$ gives
\begin{equation}\label{eq:euler}
 \int_\delta^1s^{2-k}E_{k,a}(st)\,ds
 =G_{k,a}(t)-\delta^3t^ke^{-a\delta^2t^2}.
\end{equation}
The left side is an $L^1$ limit of Riemann sums of elements of
$\mathcal C$. Here
$E_{k,a}(st)=s^kE_{k,as^2}(t)$, so every summand belongs to
$\mathcal C$. The remainder satisfies the exact identity
\[
 \norm{\delta^3t^ke^{-a\delta^2t^2}}_1
 =\delta^{2-k}\norm{G_{k,a}}_1\longrightarrow0.
\]
It follows that $G_{0,a},G_{1,a}\in\mathcal C$ for every $a>0$.

For $h>0$ and an integer $j\ge0$, consider
\[
 t^ke^{-at^2}\left(\frac{1-e^{-ht^2}}{h}\right)^j
 =h^{-j}\sum_{\ell=0}^j(-1)^\ell\binom j\ell
      t^ke^{-(a+\ell h)t^2}.
\]
The right side lies in $\mathcal C$. Since
$0\le(1-e^{-ht^2})/h\le t^2$ and this quotient tends to $t^2$,
dominated convergence, with bound $|t|^{k+2j}e^{-at^2}$, gives the
$L^1$ limit $t^{k+2j}e^{-at^2}$. Every nonnegative degree is of the
form $k+2j$ with $k=0$ or $1$.
\end{proof}

\begin{statementframe}
\begin{lemma}\label{lem:density}
The function $\chi=\1_{(0,1)}$ is an $L^1(\R)$ limit of functions
$P(t)\varphi(t)$ with $P$ a real polynomial.
\end{lemma}
\end{statementframe}
\begin{proof}
Let $d\gamma(t)=\varphi(t)\,dt$. We first show that real polynomials
are dense in $L^2(\gamma;\R)$. If $h$ is orthogonal to all polynomials,
the function
\[
 F(z)=\int_\R h(t)e^{zt}\,d\gamma(t),\qquad z\in\C,
\]
is entire. Indeed, on every bounded set of $z$ its integrand and all
its derivatives are dominated by integrable functions, by
Cauchy--Schwarz and the Gaussian exponential moments. All derivatives
of $F$ at zero vanish by orthogonality, so $F$ is identically zero.
In particular the Fourier transform of the finite signed measure
$h\,d\gamma$ vanishes. Uniqueness of the Fourier transform for finite
measures gives $h=0$ almost everywhere. The orthogonal-complement
criterion proves density.

Now $\chi/\varphi\in L^2(\gamma)$, since
\[
 \int_\R\left|\frac{\chi}{\varphi}\right|^2\,d\gamma
 =\int_0^1\frac{dt}{\varphi(t)}<\infty.
\]
Choose polynomials $P_j\to\chi/\varphi$ in $L^2(\gamma)$. As
$\gamma(\R)=1$, Cauchy--Schwarz gives
\[
 \norm{P_j\varphi-\chi}_{L^1(dt)}
 \le\norm{P_j-\chi/\varphi}_{L^2(\gamma)}\longrightarrow0.
\]
\end{proof}

\begin{statementframe}
\begin{proposition}\label{prop:interval}
If \eqref{eq:flowaverage} converges almost everywhere for each positive
dilation of $E$ and $O$, then
\[
 \frac1N\int_0^N f(U_ty)g(V_ty)\,dt
\]
converges almost everywhere for bounded measurable $f,g$.
\end{proposition}
\end{statementframe}
\begin{proof}
Apply Lemma~\ref{lem:closure} to $H_N(y,t)=f(U_{Nt}y)g(V_{Nt}y)$.
Its closed subspace contains all the Euler dilations by hypothesis,
all polynomial Gaussian kernels by Lemma~\ref{lem:gaussian}, and
then $\chi$ by Lemma~\ref{lem:density}. Finally, the change of variables
$s=Nt$ identifies the pairing with $\chi$ with the displayed time
average. Only countably many kernels are used in each limiting
argument, so no uncountable intersection of sets of full measure is
required.
\end{proof}

\section{Suspension and exact reconstruction}

\begin{statementframe}
\begin{lemma}\label{lem:suspension}
On $Y=X\times[0,1)^2$ with $\nu=\mu\otimes du\otimes dv$, the maps
\begin{align*}
 U_t(x,u,v)&=(T^{\lfloor u+t\rfloor}x,\{u+t\},v),\\
 V_t(x,u,v)&=(S^{\lfloor v+t\rfloor}x,u,\{v+t\})
\end{align*}
define commuting, jointly measurable, measure-preserving $\R$-actions.
Here integer powers use the inverse transformations at negative times.
\end{lemma}
\end{statementframe}
\begin{proof}
The floor identity
\[
 \lfloor u+s+t\rfloor
 =\lfloor u+s\rfloor+\lfloor\{u+s\}+t\rfloor
\]
gives the group law for $U$, and the same argument gives it for $V$.
The two maps act on different phase coordinates, and their actions on
$X$ commute because all integer powers of $T$ and $S$ commute.
Joint measurability follows by partitioning the time-phase plane into
the countably many measurable sets where its floor is a fixed integer.

For fixed $t$, write $t=q+a$ with $q\in\mathbb Z$ and $0\le a<1$.
Partition the $u$-interval into $[0,1-a)$ and $[1-a,1)$. On these two
pieces $U_t$ applies respectively $T^q$ and $T^{q+1}$ to $x$ and
translates $u$ onto $[a,1)$ and $[0,a)$. Invariance of $\mu$ and
translation invariance of Lebesgue measure prove preservation of the
product measure. The proof for $V_t$ is identical.
\end{proof}

Lift $f,g$ to $Y$ by ignoring the phase coordinates. Their interval-kernel
average is
\[
 B_N(x,u,v)=\frac1N\int_0^N
 f(T^{\lfloor t+u\rfloor}x)g(S^{\lfloor t+v\rfloor}x)\,dt.
\]

\begin{statementframe}
\begin{lemma}\label{lem:phase}
Let $w(u,v)=3-6|u-v|$. For every bounded $f,g$, every $x$, and $N\ge1$,
\begin{equation}\label{eq:phase}
\begin{split}
 &\int_0^1\!\int_0^1 w(u,v)B_N(x,u,v)\,dv\,du\\
 &\qquad=A_N(f,g)(x)+\frac{f(T^Nx)g(S^Nx)-f(x)g(x)}{2N}.
\end{split}
\end{equation}
\end{lemma}
\end{statementframe}
\begin{proof}
For fixed $n$, abbreviate
\begin{align*}
 a&=f(T^nx)g(S^nx),& b&=f(T^{n+1}x)g(S^nx),\\
 c&=f(T^nx)g(S^{n+1}x),& d&=f(T^{n+1}x)g(S^{n+1}x).
\end{align*}
Splitting the time interval $[n,n+1]$ at $n+1-u$ and $n+1-v$ gives
the cell integral
\[
 (1-\max(u,v))a+(u-v)_+b+(v-u)_+c+\min(u,v)d.
\]
The weighted square integrals of the four coefficients are
\begin{align*}
 \int_{[0,1]^2}w(1-\max(u,v))&=\tfrac12,&
 \int_{[0,1]^2}w(u-v)_+&=0,\\
 \int_{[0,1]^2}w(v-u)_+&=0,&
 \int_{[0,1]^2}w\min(u,v)&=\tfrac12.
\end{align*}
For example, on $u\ge v$ use $q=u-v$; the second integral is
$\int_0^1(3-6q)q(1-q)\,dq=0$. The first is
$\int_0^1(3-6q)(1-q)^2\,dq=1/2$, after adding the two triangles.
Symmetry proves the other two identities. Thus each cell contributes
$(a+d)/2$. Summing in $n=0,\ldots,N-1$ and dividing by $N$ telescopes
to \eqref{eq:phase}.
\end{proof}

\begin{statementframe}
\begin{proposition}\label{prop:bounded}
The averages $A_N(f,g)$ converge almost everywhere whenever
$f,g\in L^\infty(\mu)$.
\end{proposition}
\end{statementframe}
\begin{proof}
Use bounded measurable representatives. Apply
Corollary~\ref{cor:euler} and Proposition~\ref{prop:interval} to the
actions of Lemma~\ref{lem:suspension}. This proves convergence of
$B_N$ for $\nu$-almost every $(x,u,v)$. By Fubini, for almost every
$x$ the convergence holds for almost every phase pair. Since
\[
 |B_N(x,u,v)|\le\norm f_\infty\norm g_\infty,
 \qquad |w(u,v)|\le3,
\]
dominated convergence gives convergence of the left side of
\eqref{eq:phase} for such $x$. The boundary term on its right side is
bounded in modulus by $\norm f_\infty\norm g_\infty/N$ and tends to
zero. Hence $A_N(f,g)(x)$ converges. The argument integrates over the
whole phase square and does not restrict a product null set to a
diagonal.
\end{proof}

\section{Extension to square-integrable functions}

\begin{statementframe}
\begin{proposition}\label{prop:l2}
Fix two measure-preserving transformations $T,S$ on a probability space.
If $A_N(u,v)$ converges almost everywhere for every pair of bounded
measurable functions $u,v$, then it converges almost everywhere for
every $f,g\in L^2(\mu)$.
\end{proposition}
\end{statementframe}
\begin{proof}
Write $f_M=f\1_{\{|f|\le M\}}$ and $g_M=g\1_{\{|g|\le M\}}$, with
$M$ a positive integer, and set
\[
 P_N^Th(x)=\frac1N\sum_{n=0}^{N-1}h(T^nx).
\]
Cauchy--Schwarz in the finite sum gives
\begin{align}\label{eq:tail}
 |A_N(f,g)-A_N(f_M,g_M)|
 &\le \bigl(P_N^T|f-f_M|^2\bigr)^{1/2}
          \bigl(P_N^S|g|^2\bigr)^{1/2}\notag\\
 &\quad+\bigl(P_N^T|f|^2\bigr)^{1/2}
          \bigl(P_N^S|g-g_M|^2\bigr)^{1/2}.
\end{align}
Let $\mathcal I_T,\mathcal I_S$ be the invariant sigma-algebras and put
\[
 F_M=\E(|f-f_M|^2\mid\mathcal I_T),\quad
 G_M=\E(|g-g_M|^2\mid\mathcal I_S),
\]
\[
 F=\E(|f|^2\mid\mathcal I_T),\qquad
 G=\E(|g|^2\mid\mathcal I_S).
\]
The pointwise ergodic theorem identifies the four limits in
\eqref{eq:tail} with $F_M,G,F,G_M$, respectively. These statements
hold simultaneously outside a null set, since the family of truncation
levels is countable. The functions $F,G$ are finite almost everywhere.

The tail integrals tend to zero, so
$\norm{F_M}_1+\norm{G_M}_1\to0$. Choose a strictly increasing sequence
$M_j$ for which this sum is at most $2^{-j}$. Tonelli's theorem gives
$\sum_j(F_{M_j}+G_{M_j})<\infty$ almost everywhere, and hence both
$F_{M_j}$ and $G_{M_j}$ tend to zero there. By \eqref{eq:tail},
\[
 \limsup_{N\to\infty}|A_N(f,g)-A_N(f_{M_j},g_{M_j})|
 \le F_{M_j}^{1/2}G^{1/2}+F^{1/2}G_{M_j}^{1/2}
 \longrightarrow0.
\]
Each truncated average converges by assumption. Its asymptotic
oscillation is therefore zero. The triangle inequality bounds the
asymptotic oscillation of $A_N(f,g)$ by twice the displayed error for
every $j$. Letting $j\to\infty$ proves the Cauchy property and therefore
convergence almost everywhere.
\end{proof}

\begin{proof}[Proof of Theorem~\ref{thm:main}]
Proposition~\ref{prop:bounded}, using the planar estimate proved in
Proposition~\ref{prop:planar}, gives convergence for bounded inputs.
Proposition~\ref{prop:l2} extends it to
$L^2\times L^2$. On a probability space, $L^\infty$ and $L^3$ are
contained in $L^2$, which proves all the stated conclusions.
\end{proof}


\appendix
\section{Traceable matrix and heat inputs}\label{app:inputs}
All nonstandard matrix and heat inputs are taken from the source of
\cite{FrozenMatrix,FrozenHeat,FrozenPreliminaries}, frozen at commit
\texttt{adc7f1241b42e322a6451854ab7e4b4c146bf78a}.
The accompanying artifact records the SHA256 hashes in
\texttt{research/public-proof/receipt.json}; the saved source files
remain unmodified. The following identifiers are literal source labels,
so a reader can locate the proofs without relying on changing theorem
numbers.
\begin{description}
\item[Matrix:] \path{lem:hessian}, \path{lem:order},
\path{lem:joint-convexity}, \path{lem:rank-change}, and
\path{lem:mixed-trace}. These supply the supported Hessian, cost
order and probability averaging, rank-change integral, and Hermitian
mixed cyclic trace estimate used here.
\item[Heat:] \texttt{lem:heat-identities},
\texttt{lem:initial-energy}, and \texttt{lem:single-edge-energy}.
These include the integrated plane-chart identities, local integrable
Gaussian domination, and the factor $M=1+\sqrt{2/\pi}$ in the initial
single-edge bound.
\item[Maximal averages:] \texttt{lem:line-maximal} in the
preliminaries supplies the strong $L^{3/2}$ estimate and discrete
Gaussian domination whenever the variance is at least $h^2$.
\end{description}

\section{Formal statement and scope}\label{app:formal}
The accompanying Lean project is pinned to Lean 4.33.1 and to the following
Mathlib commit:
\begin{quote}
\texttt{0df444a360eaa60ab8c11dca51a86af692955474}.
\end{quote}
The qualitative theorem for arbitrary probability spaces, commuting
invertible bimeasurable measure-preserving maps, and complex $L^2$
inputs is \path{commuting_pointwise_l2} in
\texttt{UnconditionalConvergence.lean}. It assumes no ergodicity and
no unproved analytic target. The analytic input is internally supplied
by \path{eulerCompactPlaneTarget_proved}.

The file \texttt{SubmissionStatements.lean} also provides
\path{commuting_measurable_inverses_l2_convergence}.
It directly quantifies the maps and their measurable left and right
inverses, and concludes that the sums over $n=0,\ldots,N-1$ and
$n=1,\ldots,N$ converge to the same complex limit almost everywhere.
Real-valued $L^2$, $L^3$, and $L^\infty$ forms have real limits.
The theorem \path{ae_same_average_limits_of_representatives}
supplies one full-measure set on which changing representatives leaves
every average, and hence every candidate limit, unchanged.

The iteration-70 project audit covers 2095 public theorem and lemma
declarations, with only \texttt{propext}, \texttt{Classical.choice},
and \texttt{Quot.sound}. This audit is an incremental compilation with
source-hash matching and a fresh full axiom check. It is distinct from
the separate fresh compilation of all required dependency sources.
The formalized planar input is the actual finite maximum over positive
natural chains inside the spatial integral, sufficient for
Theorem~\ref{thm:main}. The artifact does not label every quantitative
all-real-scale assertion in this manuscript as independently formalized.
Mathematical peer review, independent dependency rebuilding, and
publication are separate statuses.

\begin{thebibliography}{9}
\bibitem{FrozenMatrix}
OpenAI, \emph{Annular variation of the triangular Hilbert transform at
the symmetric point}, October 5, 2026, dimension-free matrix
inequalities, frozen source:
\href{https://github.com/openai/math/blob/adc7f1241b42e322a6451854ab7e4b4c146bf78a/preprints/Annular-variation-of-the-triangular-Hilbert-transform-at-the-symmetric-point-October-5-2026/build/sections/matrix.tex}{\texttt{matrix.tex}}.
\bibitem{FrozenHeat}
OpenAI, same manuscript and commit, Gaussian heat dissipation:
\href{https://github.com/openai/math/blob/adc7f1241b42e322a6451854ab7e4b4c146bf78a/preprints/Annular-variation-of-the-triangular-Hilbert-transform-at-the-symmetric-point-October-5-2026/build/sections/heat.tex}{\texttt{heat.tex}}.
\bibitem{FrozenPreliminaries}
OpenAI, same manuscript and commit, truncations, coordinates, and
maximal averages:
\href{https://github.com/openai/math/blob/adc7f1241b42e322a6451854ab7e4b4c146bf78a/preprints/Annular-variation-of-the-triangular-Hilbert-transform-at-the-symmetric-point-October-5-2026/build/sections/preliminaries.tex}{\texttt{preliminaries.tex}}.
\end{thebibliography}

\end{document}
